% Supplement to pa_bridge.tex.  No new packages are required.
\section{Alternative implementations of the final AX jump}
\label{sec:closing-alternatives}

This section records a finite implementation experiment.  All state numbers
in this section are the imported native labels, before dense relabelling.
The closing instruction of the AX emitter is \texttt{breakTo 1}.  Its reader
transition on its final address bit is
\[
       (q,1)\longmapsto (0,\mathrm L,302).
\]
The seven-head source construction changes the emitted formulas, but leaves
the final guard and cleanup macros unchanged.  The original source proof,
\texttt{PA196/Source/AXGate.lean}, in the local claim named
\texttt{fifth}, establishes that logical registers $3,4,5,6,7$ are all zero
immediately before this instruction.  This assertion is an arithmetic
invariant, independent of any assumption about the consistency of PA.
The correspondence of these fields with register numbers is in
\texttt{PA196/Source/AXGateSteps.lean}.

\paragraph{Known-zero conditional decrement.}
For a source address $p$, both \texttt{breakTo 1} and the zero branch of
\texttt{dec}$\,r$ cut the address at its innermost zero edge and continue
past that subtree.  A zero decrement changes no register.  Consequently,
at this particular source site, the source instructions have identical
successors and register effects whenever register $r$ is zero.  Existing
native entries give the four alternatives
\[
\begin{array}{c|cccc}
r&3&4&5&6\\ \hline
\text{closing action}&1\mathrm R17&1\mathrm R144&
  1\mathrm R45&1\mathrm R21 .
\end{array}
\]
These four alternatives were already supplied in the earlier
\texttt{pa\_tm\_extended\_research\_2026-11-13} archive, under
\texttt{auxiliary196/closing\_variants.md}.
For each modified physical table, all source instruction contracts must
still be checked.  Only the closing contract is restricted to the zero
branch of the selected register.  It would be incorrect to impose this
restriction at every occurrence of a conditional decrement, or to validate
the modified table against the old closing row while leaving the source
interpretation implicit.

\paragraph{An unconditional padding detour.}
There is also an alternative which needs no register invariant:
\[
                  (q,1)\longmapsto(1,\mathrm R,113).
\]
It requires only that the cell immediately following the completed address
contains zero.  This follows from the explicit padding bound in the
source-frame contract.  The two relevant unchanged native rows are
\[
 (113,0)\longmapsto(0,\mathrm L,302),\qquad
 (302,1)\longmapsto(0,\mathrm L,302).
\]
Let the head initially be at address cell $i$, containing $1$, and suppose
cell $i+1$ contains $0$.  The new transition leaves the $1$ intact and moves
to $i+1$ in state $113$.  The next transition leaves its $0$ intact and
returns to $i$ in state $302$.  The third transition changes the $1$ at $i$
to $0$ and moves to $i-1$ in state $302$.  This is precisely the complete
configuration reached by the original closing transition after one step.
Thus the detour costs two literal transitions and preserves every tape
cell, the final head position, and the final control state.  The proof is
uniform in all register values and in every suffix of the tape.  It is a
finite calculation that PA verifies directly.

The necessary padding condition must be checked for the actual grammar.
For the searched AX grammars the AX-subtree height is at most $38$, its
fixed exterior prefix has length $8$, and source frames provide
$48-|p|+u$ padding zeros, with $u\geq0$.  Hence even the longest AX address
has at least two following zeros.  A future grammar search changing these
bounds must re-establish the condition instead of assuming it.

\paragraph{Three alternatives using the preceding address bit.}
The final AX leaf is the rightmost leaf of its grammar tree.  Its path
within that tree therefore consists entirely of ones.  The fixed exterior
prefix also ends in one.  Consequently, immediately before the closing
leaf is executed, both its current address bit and its preceding address
bit are one.  No register-value assumption is used here.
The fixed 134-state portion of the table contains states $88$, $103$ and
$301$, and its rows include
\[
 (88,1)\longmapsto(1,\mathrm R,103),\qquad
 (301,1)\longmapsto(1,\mathrm R,103),\qquad
 (103,1)\longmapsto(0,\mathrm L,302).
\]
Replacing the closing action by either $1\mathrm L88$ or
$1\mathrm L301$ first moves to the preceding one, moves back to the
closing one in state $103$, and then erases that one and moves left in
state $302$.  These three transitions reach exactly the configuration
reached by the original closing action in one transition.  Every other
tape cell is unchanged.

The replacement $0\mathrm L103$ also works.  After its first step,
compare it with the original $0\mathrm L302$: both runs have erased the
closing one and have their head on the preceding one.  Their controls
are $103$ and $302$, respectively.  Since the two controls have the same
row on symbol one, both runs reach the same complete configuration after
their second step.  This is a bounded convergence argument with two
steps on each side, rather than a change to the source instruction.

These arguments require the listed fixed rows to remain unchanged.
They remain valid if compatible reader folds elsewhere add finite
stuttering, provided the final physical contract is checked for the
folded table.  Evaluating these three actions and the additional
register-6-zero action $1\mathrm L77$ on all 98 archived tied grammars
adds 392 cases; none has a matching-based reader score below 196.

\paragraph{Exact scope of the finite experiment.}
The script \texttt{closing\_ties.py} independently reconstructs the 98
archived tied grammars and four distinct current gate-search grammars.
For all 102 grammars it evaluates seven choices of closing action: the
original jump, the four known-zero alternatives, the unconditional
$1\mathrm R113$ detour, and the separately checked register-6-zero route
$1\mathrm R103$.  The latter route is retained only as a checked search
case; the simple three-transition argument above applies to $113$.
There are 714 grammar/action cases.  None has an effective score below
196 under the specified matching-based reader-fold operation.  This is
only an exhaustive statement about that finite family and that operation,
not a lower bound for the machine problem.

For a selected minimum in each mode, the full 4780 instruction contracts
pass, with maximum abstract execution length 506.  These contracts use
an arbitrary-natural padding parameter and the complete arbitrary
register-file suffix, restricted only at the selected closing instruction
when its proof uses a zero register.  A separate recursive comparison
finds no additional exact-row reuse with existing fixed reader states in
the 98 archived grammars under these seven closing modes.  The JSON
receipts identify the exact selected grammars and folds.  Abstract
execution lengths count symbolic checker steps, not literal TM steps.
No complete new Lean endpoint is asserted for these experiments.
